% Part: normal-modal-logic
% Chapter: completeness
% Section: complete-consistent-sets

\documentclass[../../../include/open-logic-section]{subfiles}

\begin{document}

\olfileid{nml}{com}{ccs}

\olsection{المجموعات $\Sigma$-المتسقة التامة}

افترض أن~$\Sigma$ مجموعة من ال!!{formula}s الجهية---وفكر فيها بوصفها
بديهيات منطق جهوي نظامي أو مبادئه المعرِّفة. تكون المجموعة~$\Gamma$
$\Sigma$-متسقة إذا وفقط إذا كان~$\Gamma \Proves/[\Sigma] \lfalse$،
أي إذا لم يوجد !!{derivation} لـ
$!A_1 \lif (!A_2 \lif \cdots (!A_n \lif \lfalse)\dots)$ من~$\Sigma$،
حيث~$!A_i \in \Gamma$ لكل~$i$. وسنبني نموذجًا ``قانونيًا'' نتخذ فيه
كل عالم نوعًا خاصًا من المجموعات $\Sigma$-المتسقة: مجموعة لا تقتصر
على الاتساق مع~$\Sigma$، بل تكون قصوى بهذا المعنى، إذ تفصل في قيمة
صدق كل !!{formula} جهية: لكل~$!A$، إما~$!A \in \Gamma$ أو
$\lnot !A \in \Gamma$:

\begin{defn}
  تكون المجموعة~$\Gamma$ \emph{$\Sigma$-متسقة تامة} إذا وفقط إذا
  كانت $\Sigma$-متسقة، وكان، لكل~$!A$، إما~$!A \in \Gamma$ أو
  $\lnot !A \in \Gamma$.
\end{defn}

للمجموعات $\Sigma$-المتسقة التامة~$\Gamma$ خصائص نافعة عدة. فهي،
أولًا، مغلقة استنتاجيًا؛ أي إنه إذا كان~$\Gamma \Proves[\Sigma] !A$،
كان~$!A \in \Gamma$. ويعني هذا، خصوصًا، أن كل مثال من
!!a{formula}~$!A \in \Sigma$ ينتمي أيضًا إلى~$\Gamma$. وفوق ذلك،
تعكس العضوية في~$\Gamma$ شروط الصدق للروابط القضوية. وسيكون هذا مهمًا
حين نعرّف ``النموذج القانوني''.

\begin{prop}\ollabel{prop:ccs-properties}
  افترض أن~$\Gamma$ مجموعة $\Sigma$-متسقة تامة. عندئذ:
  \begin{enumerate}
  \item \ollabel{prop:ccs-closed}%
    $\Gamma$ مغلقة استنتاجيًا في~$\Sigma$.
  \item \ollabel{prop:ccs-sigma}%
    $\Sigma \subseteq \Gamma$.
  \tagitem{prvFalse}{\ollabel{prop:ccs-lfalse}%
    $\lfalse \notin \Gamma$}{}
  \tagitem{prvTrue}{\ollabel{prop:ccs-ltrue}%
    $\ltrue \in \Gamma$}{}
  \item \ollabel{prop:ccs-lnot}%
    $\lnot!A \in \Gamma$ إذا وفقط إذا كان~$!A \notin \Gamma$.
  \tagitem{prvAnd}{\ollabel{prop:ccs-land}%
    $!A \land !B \in \Gamma$ إذا وفقط إذا كان~$!A \in \Gamma$ و
    $!B \in \Gamma$}{}
  \tagitem{prvOr}{\ollabel{prop:ccs-lor}%
    $!A \lor !B \in \Gamma$ إذا وفقط إذا كان~$!A \in \Gamma$ أو
    $!B \in \Gamma$}{}
  \tagitem{prvIf}{\ollabel{prop:ccs-lif}%
    $!A \lif !B \in \Gamma$ إذا وفقط إذا كان~$!A \notin \Gamma$ أو
    $!B \in \Gamma$}{}
  \tagitem{prvIff}{\ollabel{prop:ccs-liff}%
    $!A \liff !B \in \Gamma$ إذا وفقط إذا كان إما~$!A \in \Gamma$
    و$!B \in \Gamma$، أو~$!A \notin \Gamma$ و$!B \notin \Gamma$}{}
  \end{enumerate}
\end{prop}

\begin{proof}
  \begin{enumerate}
  \item افترض أن~$\Gamma \Proves[\Sigma] !A$ لكن~$!A \notin \Gamma$.
    ولما كانت~$\Gamma$ $\Sigma$-متسقة تامة، كان~$\lnot!A \in \Gamma$.
    وهذا يجعل~$\Gamma$ غير متسقة، لأن
    $!A, \lnot !A \Proves[\Sigma] \lfalse$.

  \item إذا كان~$!A \in \Sigma$، فإن~$\Gamma \Proves[\Sigma] !A$،
    ومن ثم~$!A \in \Gamma$ بالإغلاق الاستنتاجي، أي بالحالة
    \olref{prop:ccs-closed}.

  \tagitem{prvFalse}{إذا كان~$\lfalse \in \Gamma$، فإن
    $\Gamma \Proves[\Sigma] \lfalse$، فتكون~$\Gamma$
    $\Sigma$-غير متسقة.}{}

  \tagitem{prvTrue}{$\Gamma \Proves[\Sigma] \ltrue$، ومن ثم
    $\ltrue \in \Gamma$ بالإغلاق الاستنتاجي، أي بالحالة
    \olref{prop:ccs-closed}.}{}

  \item إذا كان~$\lnot!A \in \Gamma$، فإن الاتساق يقتضي
    $!A \notin \Gamma$؛ وإذا كان~$!A \notin \Gamma$، فإن
    $\lnot !A \in \Gamma$ لأن~$\Gamma$ $\Sigma$-متسقة تامة.

  \tagitem{prvAnd}{\iftag{probAnd}{تمرين.}{افترض أن
      $!A \land !B \in \Gamma$. ولما كانت
      $(!A \land !B) \lif !A$ مثال إحلال لقضية صادقة صدقًا تحليليًا،
      كان~$!A \in \Gamma$ بالإغلاق الاستنتاجي، أي بالحالة
      \olref{prop:ccs-closed}. وكذلك~$!B \in \Gamma$. ومن ناحية أخرى،
      افترض أن~$!A \in \Gamma$ و$!B \in \Gamma$. عندئذ يقتضي الإغلاق
      الاستنتاجي أن~$(!A \land !B) \in \Gamma$، لأن
      $!A \lif (!B \lif (!A \land !B))$ مثال إحلال لقضية صادقة صدقًا
      تحليليًا.}}{}

  \tagitem{prvOr}{\iftag{probOr}{تمرين.}{افترض أن
      $!A \lor !B \in \Gamma$، وأن~$!A \notin \Gamma$ و
      $!B \notin \Gamma$. ولما كانت~$\Gamma$ $\Sigma$-متسقة تامة،
      كان~$\lnot!A \in \Gamma$ و$\lnot !B \in \Gamma$. ومن ثم
      $\lnot(!A \lor !B) \in \Gamma$، لأن
      $\lnot !A \lif (\lnot !B \lif \lnot (!A \lor !B))$ مثال إحلال
      لقضية صادقة صدقًا تحليليًا. وهذا يعني أن~$\Gamma$
      $\Sigma$-غير متسقة، وهو تناقض. والعكس مباشر: إذا كان
      $!A \in \Gamma$ أو~$!B \in \Gamma$، فإن الإغلاق الاستنتاجي
      يقتضي~$!A \lor !B \in \Gamma$.}}{}

  \tagitem{prvIf}{\iftag{probIf}{تمرين.}{افترض أن
      $!A \lif !B \in \Gamma$ و$!A \in \Gamma$؛ عندئذ
      $\Gamma \Proves[\Sigma] !B$، ومن ثم~$!B \in \Gamma$ بالإغلاق
      الاستنتاجي. وبالعكس، إذا كان~$!A \lif !B \notin \Gamma$، فبما
      أن~$\Gamma$ $\Sigma$-متسقة تامة، يكون
      $\lnot (!A \lif !B) \in \Gamma$. ولما كانت
      $\lnot(!A \lif !B) \lif !A$ مثال إحلال لقضية صادقة صدقًا
      تحليليًا، كان~$!A \in \Gamma$ بالإغلاق الاستنتاجي. ولما كانت
      $\lnot(!A \lif !B) \lif \lnot !B$ مثالًا كذلك، كان
      $\lnot !B \in \Gamma$. ومن ثم~$!B \notin \Gamma$ لاتساق
      $\Gamma$ مع~$\Sigma$.}}{}

  \tagitem{prvIff}{\iftag{probIff}{تمرين.}{افترض أن
      $!A \liff !B \in \Gamma$. إذا كان~$!A \in \Gamma$، كان
      $!B \in \Gamma$، لأن~$(!A \liff !B) \lif (!A \lif !B)$ مثال
      إحلال لقضية صادقة صدقًا تحليليًا. وبالمثل، إذا كان
      $!B \in \Gamma$، كان~$!A \in \Gamma$. ومن ثم فإما أن يكون
      $!A \in \Gamma$ و$!B \in \Gamma$، أو لا يكون
      $!A \in \Gamma$ ولا~$!B \in \Gamma$.

      وبالعكس، افترض أن~$!A \liff !B \notin \Gamma$. ولما كانت
      $\Gamma$ $\Sigma$-متسقة تامة، كان
      $\lnot (!A \liff !B) \in \Gamma$. ولما كانت
      $\lnot(!A \liff !B) \lif (!A \lif \lnot !B)$ مثال إحلال لقضية
      صادقة صدقًا تحليليًا، فإذا كان~$!A \in \Gamma$، كان
      $\lnot !B \in \Gamma$، ومن ثم~$!B \notin \Gamma$ لاتساق
      $\Gamma$ مع~$\Sigma$. وبالمثل، إذا كان~$!B \in \Gamma$، كان
      $!A \notin \Gamma$. ومن ثم لا يكون~$!A \in \Gamma$ و
      $!B \in \Gamma$ معًا، ولا يكون~$!A \notin \Gamma$ و
      $!B \notin \Gamma$ معًا.}}{}
  \end{enumerate}
\end{proof}

\begin{probtag}{probAnd,probOr,probIf,probIff}
  أكمل برهان~\olref[nml][com][ccs]{prop:ccs-properties}.
\end{probtag}

\end{document}
